\originalchapter{有限状态马尔可夫链}
\sourcelecture{32}

\originalsection{状态、历史与转移矩阵}
独立试验每次重新开始, 下一次的分布不受过去影响. 另一类系统有持续的状态: 机器是否运行、队列所处级别、天气类别. 马尔可夫模型允许下一步依赖现在, 同时假设现在的状态已包含预测下一步所需的全部过去信息.
\begin{definition}
设状态空间 $S=\{0,\ldots,d-1\}$ 有限. 序列 $X_0,X_1,\ldots$ 是时间齐次马尔可夫链, 若存在矩阵 $P=(p_{ij})$, 使每个概率为正的历史满足
\[\Prob(X_{n+1}=j\mid X_0=i_0,\ldots,X_n=i)=p_{ij}.\]
矩阵的每个元素非负, 每行和为1, 称为转移矩阵. 初始分布记为行向量 $\alpha=(\alpha_0,\ldots,\alpha_{d-1})$.
\end{definition}
齐次是指 $p_{ij}$ 不依赖时刻 $n$. 条件概率只在条件事件概率为正时使用; 不必给不可能的历史赋予额外意义. 给定初始分布和转移矩阵, 一条有限路径的概率由乘法公式确定:
\[\Prob(X_0=i_0,\ldots,X_n=i_n)=\alpha_{i_0}\prod_{r=0}^{n-1}p_{i_ri_{r+1}}.\]
这是把抽象定义转换为可计算模型的第一步.

\begin{example}
某设备每天记录为正常状态0或维护状态1. 正常后下一天进入维护的概率为 $1/4$, 维护后下一天恢复正常的概率为 $1/2$. 则
\[P=\begin{pmatrix}3/4&1/4\\1/2&1/2\end{pmatrix}.\]
从正常开始, 连续两天仍正常的路径概率为 $(3/4)^2=9/16$. 两天后正常的概率还要加上先维护再恢复的路径, 得 $9/16+(1/4)(1/2)=11/16$. 单一路径概率与末端事件概率不是同一个量.
\end{example}
\begin{center}
\begin{tikzpicture}[>=Stealth,every node/.style={font=\small}]
\node[draw,circle] (a) at (0,0) {正常};
\node[draw,circle] (b) at (4,0) {维护};
\draw[->] (a) edge[bend left=20] node[above] {$1/4$} (b);
\draw[->] (b) edge[bend left=20] node[below] {$1/2$} (a);
\draw[->] (a) edge[loop left] node[left] {$3/4$} (a);
\draw[->] (b) edge[loop right] node[right] {$1/2$} (b);
\end{tikzpicture}
\end{center}

\begin{proposition}[Chapman--Kolmogorov公式]
写 $p_{ij}^{(n)}=\Prob(X_n=j\mid X_0=i)$, 其中从 $i$ 出发的链按相同转移规则定义. 则 $P^{(0)}=I$, $P^{(n)}=P^n$, 且
\[p_{ij}^{(m+n)}=\sum_{k\in S}p_{ik}^{(m)}p_{kj}^{(n)}.\]
时刻 $n$ 的分布是 $\alpha P^n$.
\end{proposition}
\begin{proof}
按中间状态 $X_m=k$ 分解事件. 路径概率的乘法式说明, 给定 $X_m=k$ 和任何此前历史, 之后 $n$ 步到 $j$ 的概率只依赖 $k$, 等于 $p_{kj}^{(n)}$. 全概率公式给出所列求和. 取 $n=1$ 归纳得到矩阵幂. 再按初始状态求和, 第 $j$ 个概率为 $\sum_i\alpha_i(P^n)_{ij}$.
\end{proof}
例如设备模型有
\[P^2=\begin{pmatrix}11/16&5/16\\5/8&3/8\end{pmatrix}.\]
行表示起点, 列表示终点. 采用行分布向量时, 必须在右侧乘 $P$, 不能把行列约定混用.

三种特殊矩阵帮助理解定义. 若所有行等于同一分布 $q$, 则 $X_1,X_2,\ldots$ 独立同分布为 $q$, 并且与 $X_0$ 独立; 只有初始分布也为 $q$ 时, 整个 $X_0,X_1,\ldots$ 才独立同分布. 若 $P=I$, 状态永不改变, 却可以有随机的起点. 若每个元素是0或1, 每行恰有一个1, 以函数 $f(i)$ 表示其所在列, 则 $X_{n+1}=f(X_n)$, 初始状态确定之后路径完全确定.

\originalsection{停留时间与状态选择}
\begin{proposition}
从状态 $i$ 出发, 令 $T=\min\{n\geq1:X_n\ne i\}$. 若 $p_{ii}<1$, 则
\[\Prob(T=k)=p_{ii}^{k-1}(1-p_{ii}),\quad k\geq1,\qquad \E T=\frac1{1-p_{ii}}.\]
若 $p_{ii}=1$, 则 $T=\infty$ 几乎处处.
\end{proposition}
\begin{proof}
$T=k$ 要求前 $k-1$ 步都留在 $i$, 第 $k$ 步离开, 路径乘法给出概率. 尾和公式给出 $\E T=\sum_{n\geq0}\Prob(T>n)=\sum_{n\geq0}p_{ii}^n$. 最后一种情况每一步离开概率为零, 可数并集仍概率零.
\end{proof}
在设备模型中, 从正常到首次维护平均等待4天, 从维护到首次正常平均等待2天. $T$ 从下一天计起, 因而不把今天额外加一次.

几何停留意味着已经在状态内待了很久, 也不会改变下一步离开的概率. 这是模型的可检验限制. 若一台机器的故障概率随累计运行时间上升, 仅用“正常/故障”两个状态往往不能描述历史影响. 可以把正常状态再细分为运行年龄级别, 或记录必要的累计变量. 扩充状态能保留所需记忆, 但任意扩充不自动保证马尔可夫性, 仍要检验条件概率是否真的只依赖新状态. 源讲义用关系状态讨论同样的建模问题; 此处以设备停留时间训练相同的判断.

\originalsection{平稳分布与长期分布}
\begin{definition}
概率行向量 $\pi$ 若满足 $\pi P=\pi$, 称为平稳分布. 若以 $\pi$ 初始化, 每个时刻的分布都为 $\pi$. 若存在整数 $r\geq1$ 使 $P^r$ 的所有元素严格为正, 称 $P$ 为本原转移矩阵. 源讲义把这个性质称作遍历性.
\end{definition}
平稳的意思是整个分布不变, 不是每条路径停在某状态. 它是特征值1的左特征向量, 并须同时满足非负与归一化条件. 若用列向量解方程, 正确方程是 $(P^{\mathsf T}-I)\pi^{\mathsf T}=0$.
\begin{example}
设备模型的平稳方程为
$\pi_0=(3/4)\pi_0+(1/2)\pi_1$, $\pi_0+\pi_1=1$.
第一式给 $\pi_0=2\pi_1$, 所以 $\pi=(2/3,1/3)$. 令 $u_n=\Prob(X_n=0)$, 全概率公式给出
\[u_{n+1}=\tfrac34u_n+\tfrac12(1-u_n)=\tfrac12+\tfrac14u_n.\]
减去固定点 $2/3$, 得 $u_{n+1}-2/3=(u_n-2/3)/4$, 因而
\[u_n=\frac23+4^{-n}\left(u_0-\frac23\right).\]
这个计算同时给出平稳分布和任意初始分布的收敛速率, 比只解线性方程多了收敛结论.
\end{example}

\begin{theorem}[有限链的平稳分布存在]
每个有限状态转移矩阵至少有一个平稳分布.
\end{theorem}
\begin{proof}
取任意概率行向量 $\alpha$, 令 $\nu_n=n^{-1}\sum_{k=0}^{n-1}\alpha P^k$. 它们的每个坐标都在 $[0,1]$, 由有限维Bolzano--Weierstrass定理可取收敛子列 $\nu_{n_l}\to\pi$. 极限仍非负且坐标和为1. 又有
$\nu_nP-\nu_n=(\alpha P^n-\alpha)/n\to0$.
矩阵乘法是连续的, 沿子列取极限得 $\pi P=\pi$.
\end{proof}
\begin{lemma}[正矩阵幂带来的收缩]
设 $Q=P^r$ 所有元素严格为正, $\eta=\min_{i,j}Q_{ij}>0$, $\theta=d\eta\leq1$. 对任意两个概率行向量 $\alpha,\beta$,
\[\|\alpha Q-\beta Q\|_1\leq(1-\theta)\|\alpha-\beta\|_1,\qquad \|v\|_1=\sum_i|v_i|.\]
\end{lemma}
\begin{proof}
若 $\theta=1$, 每行和为1迫使 $Q_{ij}=1/d$, 两个分布乘 $Q$ 后都成为均匀分布. 若 $\theta<1$, 令 $U$ 每个元素为 $1/d$, $R_{ij}=(Q_{ij}-\eta)/(1-\theta)$. $R$ 非负, 行和1, 且 $Q=\theta U+(1-\theta)R$. 因概率向量之差的坐标和为零, $(\alpha-\beta)U=0$. 对任何行向量 $v$,
\[\|vR\|_1\leq\sum_j\sum_i|v_i|R_{ij}=\sum_i|v_i|.\]
代入 $v=\alpha-\beta$ 即得结论. 同一计算也说明乘任意转移矩阵不会增大此距离.
\end{proof}
\begin{theorem}[本原有限链的收敛]
若 $P$ 本原, 则平稳分布唯一, 且对任意初始分布 $\alpha$,
\[\alpha P^n\to\pi,\qquad (P^n)_{ij}\to\pi_j.\]
若 $P^r$ 的最小元素为 $\eta$, 则
$\|\alpha P^n-\pi\|_1\leq2(1-d\eta)^{\lfloor n/r\rfloor}$.
\end{theorem}
\begin{proof}
先由存在定理取平稳分布 $\pi$. 若 $\rho$ 也是平稳分布, 对 $Q^k=P^{rk}$ 反复收缩得到
$\|\rho-\pi\|_1\leq(1-d\eta)^k\|\rho-\pi\|_1$, 令 $k\to\infty$ 得唯一性. 写 $n=kr+s$, $0\leq s<r$. 先乘 $Q^k$ 收缩 $k$ 次, 再乘 $P^s$ 不增加距离, 得所列界, 初始距离至多2. 取 $\alpha$ 为状态 $i$ 的点质量就得到逐行收敛.
\end{proof}

\originalsection{不可约、周期与收敛条件}
\begin{definition}
若对每对 $i,j$ 存在 $n\geq0$ 使 $(P^n)_{ij}>0$, 称链不可约. 状态 $i$ 的周期是
$\gcd\{n\geq1:(P^n)_{ii}>0\}$. 周期为1称非周期. 在不可约链中这些集合非空, 且所有状态周期相同.
\end{definition}
最后一条也需要理由. 若有 $i$ 到 $j$ 的正概率路径长 $a$, $j$ 到 $i$ 的路径长 $b$, 则 $a+b$ 是 $i$ 的返回长度; 对任一 $j$ 的返回长度 $n$, $a+n+b$ 也是 $i$ 的返回长度. 所以 $i$ 的周期同时整除 $a+b$ 与 $a+n+b$, 从而整除所有 $n$, 即整除 $j$ 的周期. 交换 $i,j$ 得二者相等.

\begin{lemma}[返回长度的整数工具]
若一组正整数的最大公因数为1, 则其中可取有限个 $a_1,\ldots,a_l$ 最大公因数为1, 且每个足够大的整数是这些数的非负整数线性组合.
\end{lemma}
\begin{proof}
逐步加入整数所得最大公因数是递减的正整数序列, 只会有限次严格下降, 最终为1, 所以可取有限组. 模 $a_1$ 看 $a_2,\ldots,a_l$ 生成的余数. 非负组合的余数集合对加法封闭, 且包含每个元素的加法逆元: 在有限循环群中, $-x=(a_1-1)x$. 因而它是子群. 最大公因数为1和B\'ezout等式说明该子群含余数1, 即含全部余数. 为每个余数 $s$ 选一个非负组合 $b_s$ 代表它. 若 $n\geq\max_s b_s$, 取其余数 $s$, 则 $n=b_s+ka_1$, $k\geq0$, 完成证明.
\end{proof}
\begin{theorem}
对有限状态链, 本原等价于不可约且非周期. 因此不可约非周期的有限链满足上一节的唯一平稳分布与分布收敛结论.
\end{theorem}
\begin{proof}
若 $P^r>0$, 则所有状态互通. 而 $P^{r+1}>0$: 对任意列 $j$, 不可约性保证存在某 $k$ 使 $p_{kj}>0$, 故 $(P^{r+1})_{ij}\geq(P^r)_{ik}p_{kj}>0$. 每个状态可在 $r$ 与 $r+1$ 步返回, 周期为1.

反过来固定状态 $o$. 由非周期性, $o$ 的返回长度最大公因数为1. 引理选出的有限返回长度都有正概率路径; 拼接这些路径说明 $o$ 可以在每个充分大的步数返回. 对每个 $i,j$, 选一条 $i$ 到 $o$ 的路径长 $a_i$ 和一条 $o$ 到 $j$ 的路径长 $b_j$, 允许 $a_o=b_o=0$. 当 $n-a_i-b_j$ 足够大时, 中间插入一次该长度的返回路径, 就得 $(P^n)_{ij}>0$. 状态对有限, 可选一个共同的充分大 $n$, 得 $P^n>0$.
\end{proof}
\begin{example}
$P=\begin{pmatrix}0&1\\1&0\end{pmatrix}$ 不可约, 周期2. 平稳分布为 $(1/2,1/2)$, 但从0出发, 偶数时刻一定在0, 奇数时刻一定在1, 分布不收敛. 相反 $P=I$ 的每个概率向量都平稳, 起点决定最终状态. 前例显示平稳存在甚至唯一不足以保证分布收敛, 后例显示不可约条件在唯一性中发挥作用.
\end{example}

\originalsection{长期时间比例与奖励平均}
平稳分布的坐标 $\pi_j$ 还可以表示单条长期路径中访问 $j$ 的比例. 这是另一个定理, 不能从 $\Prob(X_n=j)\to\pi_j$ 直接推出, 因为各时刻的指示变量通常依赖.
\begin{theorem}[有限不可约链的时间平均]
若有限链不可约, 对任意初始分布和任意函数 $f:S\to\R$, 几乎处处有
\[\frac1n\sum_{k=0}^{n-1}f(X_k)\longrightarrow\sum_{j\in S}\pi_jf(j),\]
其中 $\pi$ 是唯一平稳分布. 此结论不需要非周期条件. 特别地, 取 $f=\ind_{\{j\}}$, 得访问状态 $j$ 的长期比例为 $\pi_j$.
\end{theorem}
\begin{proof}
固定状态 $o$. 不可约性和状态有限意味着存在 $L\geq1$, $c>0$, 从任意状态在随后 $L$ 步内访问 $o$ 的概率至少 $c$: 为每个状态选正概率路径, 取共同长度上界与路径概率的正下界. 在 $o$ 出发时要考虑正时间返回, 可以先走一步, 再使用此界, 把 $L$ 增大1即可. 每个长度 $L$ 的区块, 给定此前历史后仍有至少 $c$ 的成功概率, 所以首次访问或返回时间 $H$ 满足
$\Prob(H>mL)\leq(1-c)^m$. 它几乎处处有限, 并且所有正整数阶矩有限; 例如按区块分层, $\E H^q\leq\sum_{m\geq0}((m+1)L)^q(1-c)^m<\infty$.

先从 $o$ 出发. 令 $\tau_0=0$, $\tau_{l+1}=\min\{n>\tau_l:X_n=o\}$, 周期长度 $D_l=\tau_{l+1}-\tau_l$, 周期奖励 $V_l=\sum_{k=\tau_l}^{\tau_{l+1}-1}f(X_k)$. 这些对 $(D_l,V_l)$ 独立同分布. 这里用到的停时马尔可夫性质可直接由路径概率证明: 将 $\{\tau_l=t\}$ 按所有在 $t$ 结束的历史分解, 对每条历史, 之后任意有限路径的条件概率均为从 $o$ 开始的同一乘积, 不依赖历史; 再对 $t$ 求和. 因访问时间有限, 这些历史的概率和为1. 递归应用即得各周期独立同分布.

有限状态上 $M=\max|f|<\infty$, $|V_l|\leq MD_l$, 因而二者有有限四阶矩. 上章已证明的强大数定律给出
\[\frac{\tau_m}{m}=\frac1m\sum_{l=0}^{m-1}D_l\to d_0=\E D_0>0,\qquad
\frac1m\sum_{l=0}^{m-1}V_l\to v_0=\E V_0.\]
令 $m(n)$ 是时刻 $n$ 前已完成的周期数, 即 $\tau_{m(n)}\leq n<\tau_{m(n)+1}$. 由于 $\tau_m/m\to d_0$, 夹逼给 $m(n)/n\to1/d_0$. 未完成周期不能破坏平均: 对每个 $\eps>0$, 指数尾使 $\sum_l\Prob(D_l>\eps l)<\infty$, Borel--Cantelli给 $D_l/l\to0$ 几乎处处. 故余下奖励绝对值至多 $MD_{m(n)}=o(n)$. 时间平均于是趋于 $v_0/d_0$.

为识别这个比值, 令 $U_j$ 是一个周期中处于 $j$ 的次数, $u_j=\E U_j$. 有 $\sum_j u_j=d_0$. 对周期内从 $i$ 到 $j$ 的转移次数取期望, 得 $u_i p_{ij}$: 写成对 $k\geq0$ 的指示和, 事件 $\{k<\tau_1,X_k=i\}$ 由时刻 $k$ 的历史决定, 条件期望中的下一步概率为 $p_{ij}$. 非负求和允许逐项取期望. 一个周期的终点与起点同为 $o$, 所以每个状态的进入次数等于离开次数, 因而 $u_j=\sum_i u_i p_{ij}$. 所以 $\pi_j=u_j/d_0$ 是平稳分布, 且 $v_0/d_0=\sum_j\pi_jf(j)$.

若初始状态不是 $o$, 到首次访问 $o$ 之前的段落长度几乎处处有限, 奖励也有限, 除以 $n$ 后趋于零. 从该访问时刻开始重复以上周期证明, 得同样结论. 最后, 若 $\rho$ 是任一平稳分布, 从 $\rho$ 初始化后, 每个时间平均的期望是 $\sum_j\rho_jf(j)$. 时间平均有界于 $M$, 其几乎处处极限的期望可交换: 有界变量的收敛可由偏差事件概率趋零与 $2M$ 的上界直接证明. 因而 $\sum_j\rho_jf(j)=\sum_j\pi_jf(j)$. 逐个取指示函数得 $\rho=\pi$, 也证明了不可约链平稳分布唯一.
\end{proof}
\begin{example}
设备模型若正常日奖励为3, 维护日奖励为 $-1$, 则长期每天平均奖励几乎处处趋于
$(2/3)3+(1/3)(-1)=5/3$. 正常日与维护日并不独立, 所以不能直接把独立同分布大数定律用于每日奖励. 上述周期分解正是恢复独立结构的方法. 对周期2的交替链, 时间比例仍为 $1/2,1/2$, 虽然单个时刻的分布始终振荡; 因而时间平均与分布收敛必须分开讨论.
\end{example}

\originalsection{习题与解答}
\begin{exercise}
两状态链的矩阵为 $P=\begin{pmatrix}1-a&a\\b&1-b\end{pmatrix}$, $0<a,b<1$. 求平稳分布与任意初始分布的 $n$ 步概率. 求每次处于0后首次到1的平均等待时间.
\end{exercise}
\begin{solution}
平稳方程给 $a\pi_0=b\pi_1$, 故 $\pi=(b/(a+b),a/(a+b))$. 设 $u_n=\Prob(X_n=0)$, 则 $u_{n+1}=b+(1-a-b)u_n$. 减去固定点得
\[u_n=\frac b{a+b}+(1-a-b)^n\left(u_0-\frac b{a+b}\right).\]
因为 $|1-a-b|<1$, 分布趋稳, 若该因子负则误差交替变号. 从0离开的等待时间为成功概率 $a$ 的几何变量, 均值 $1/a$.
\end{solution}
\begin{exercise}
三状态链有 $P=\begin{pmatrix}0&1&0\\0&0&1\\1/2&0&1/2\end{pmatrix}$. 判断不可约与周期, 求平稳分布, 求状态2的长期访问比例.
\end{exercise}
\begin{solution}
路径 $0\to1\to2\to0$ 使所有状态互通, 所以不可约. 状态2有自环, 返回长度集合含1, 周期为1, 其他状态周期也为1. 平稳方程为 $\pi_0=\pi_2/2$, $\pi_1=\pi_0$, $\pi_2=\pi_1+\pi_2/2$. 归一化给 $\pi=(1/4,1/4,1/2)$. 不可约非周期有限链分布收敛到它; 独立于该分布收敛陈述, 时间平均定理保证单条路径访问2的比例几乎处处趋于 $1/2$.
\end{solution}
\begin{exercise}
设两状态转移矩阵为单位矩阵, 初始状态各以概率 $1/2$ 选取. 每个时刻的分布都是 $(1/2,1/2)$. 判断路径中状态0的访问比例是否趋于 $1/2$.
\end{exercise}
\begin{solution}
路径始终等于初始状态. 状态0的访问比例对所有 $n$ 都是 $\ind_{\{X_0=0\}}$, 因而以相等概率为1或0, 永不等于 $1/2$. 此链可约. 平稳初始化可以保证边缘分布不变, 却不能替代时间平均定理所需的不可约性.
\end{solution}
